\section{Every smooth moving centre of the original Gaussian family}
\label{sec:gt-moving-gaussian}

Keep all original parameters and fields of
Section~\ref{sec:axis-gaussian-concentration}:
$\nu,T,\tau_*,L_*>0$, real $U_*,H_*,\beta$,
$\tau=T-t$, $\sigma=\tau/\tau_*$, $L=L_*\sigma^\beta$,
$k=L^{-2}$, $A=U_*\sigma^{\beta-1}$ and $B=H_*\sigma^{-1}$.
The spatial coordinates in this section are denoted
$y=(y_1,y_2,y_3)$ before translation; thus $s=y_1^2+y_2^2$,
$z=y_3$, $\rho=|y|^2$, and $G=e^{-k|y|^2}$ in the unchanged
Gaussian formulas. This is only a naming of the input of the
translation, with its inverse given below.

Let $d:[0,T)\to\mathbb R^3$ be any smooth curve, with a smooth
one-sided extension at $0$. No bound on $d$, $d'$, or their
preterminal growth is imposed. In the original physical coordinates
$x\in\mathbb R^3$ define
\begin{align}
 u_d(x,t)&=u(x-d(t),t),&
 p_d(x,t)&=p(x-d(t),t),\notag\\
 f_d(x,t)&=f(x-d(t),t)
          -\sum_{j=1}^3d'_j(t)\partial_{y_j}u(x-d(t),t).
 \label{gt:translation}
\end{align}
The differential of $y=x-d(t)$ with respect to $x$ is $I_3$,
and its time derivative is $-d'(t)$. Hence
\[
 \partial_tu_d=(u_t-d'\cdot\nabla_yu)(x-d,t),\quad
 (u_d\cdot\nabla_x)u_d=((u\cdot\nabla_y)u)(x-d,t),\quad
 \Delta_xu_d=(\Delta_yu)(x-d,t).
\]
The pressure gradient and divergence translate in the same way.
Substitution of the already proved full Gaussian momentum equation
therefore proves
\[
 \partial_tu_d+(u_d\cdot\nabla_x)u_d+\nabla_xp_d
       =\nu\Delta_xu_d+f_d,\qquad \operatorname{div}_xu_d=0.
\]
The inverse on this triple is
\[
 u(y,t)=u_d(y+d,t),\quad p(y,t)=p_d(y+d,t),\quad
 f(y,t)=f_d(y+d,t)+d'\cdot\nabla_xu_d(y+d,t).
\]
Substitution in both orders gives the identity on all three
components. Thus the full force correction, and not only the
velocity coordinates, is part of this exact map.
Its spatial determinant is one, so the original energy and
dissipation formulas are unchanged:
\[
 \int_{\mathbb R^3}|u_d|^2\,dx=C_E\sigma^{5\beta-2},\qquad
 \int_{\mathbb R^3}|\nabla_xu_d|^2\,dx=C_D\sigma^{3\beta-2}.
\]
The constants $C_E,C_D$ are exactly those in
\eqref{eq:ag-energy-exact}. Each translated field is Schwartz at
every preterminal time. Indeed translation and differentiation
preserve that class, and $d'(t)$ is a finite vector at each such
time. This statement does not infer uniform endpoint bounds.

\subsection{Exact vorticity and force-curl translation}

Write $\omega=\operatorname{curl}_y u$ and
$C=\operatorname{curl}_y f$ for the full Gaussian expressions
\eqref{eq:ag-vorticity} and \eqref{eq:ag-curl-force}.
Spatial derivatives commute with translation, and $d'(t)$ is
independent of space. The full identities are consequently
\begin{equation}
 \operatorname{curl}_xu_d(x,t)=\omega(x-d,t),\qquad
 \operatorname{curl}_xf_d(x,t)
   =\bigl(C-d'\cdot\nabla_y\omega\bigr)(x-d,t).
 \label{gt:curl}
\end{equation}
For example, the $i$th component of the second equality is
$\sum_{j,\ell}\epsilon_{ij\ell}\partial_jf_\ell
 -\sum_{m,j,\ell}d'_m\epsilon_{ij\ell}
           \partial_j\partial_m u_\ell$ at $y=x-d$.
Commuting the two smooth spatial derivatives gives
$C_i-\sum_m d'_m\partial_m\omega_i$, which proves each sign
and every component of \eqref{gt:curl}.

\begin{theorem}[Translation cannot regularize the force of this family]
For every nonzero amplitude pair $(U_*,H_*)$ and every real
$\beta$, the translated velocity $u_d$ cannot satisfy the same
Navier--Stokes equation with a force whose first spatial
derivatives are uniformly bounded on $\mathbb R^3\times[0,T)$,
for any choice of pressure. This conclusion holds for every
smooth centre curve $d$ specified above.
\end{theorem}
\begin{proof}
First suppose $H_*\ne0$. The third vorticity component is
\[
 \omega_3(y,t)=2B e^{-k|y|^2}
                       (1-k(y_1^2+y_2^2)).
\]
It is even in each Cartesian coordinate, so
$\nabla_y\omega_3(0,t)=0$. The full poloidal vorticity has
zero third component, including when $U_*\ne0$.
Evaluating \eqref{gt:curl} at the actual moving centre gives
\begin{equation}
 (\operatorname{curl}_xf_d)_3(d(t),t)
  =\frac{2H_*}{\tau_*}\sigma^{-2}
      +\frac{20\nu H_*}{L_*^2}\sigma^{-1-2\beta}.
 \label{gt:swirl}
\end{equation}
Both terms have the sign of $H_*$. Its absolute value is at
least $2|H_*|\sigma^{-2}/\tau_*$ and tends to infinity,
irrespective of the centre velocity.

It remains to consider $H_*=0$ and $U_*\ne0$.
The complete vorticity then has the form
\[
 \omega(y,t)=(-y_2\xi,y_1\xi,0),\qquad
 \xi=AkG(10-4k|y|^2).
\]
The scalar $\xi$ is even under $y\mapsto-y$, and hence
$\omega(-y,t)=-\omega(y,t)$. Differentiating this identity
with respect to $y_j$ gives
$\partial_j\omega(-y,t)=\partial_j\omega(y,t)$.
Thus the full term $d'\cdot\nabla_y\omega$ in
\eqref{gt:curl} has identical values at opposite points,
for each actual vector $d'(t)$.

For every $c>0$ let $y_c(t)=(L\sqrt c,0,0)$.
At both $y_c$ and $-y_c$ the nonlinear part
\[
 zA^2k^2G^2(16k|y|^2-56)
\]
of $E_\xi$ vanishes exactly because $z=0$.
The remaining scalar value is
$Ak e^{-c}[P_\beta(c)/\tau+\nu kQ(c)]$.
Since $R_h=0$ when $H_*=0$, the force-curl components at
these two points are $(0,\pm L\sqrt c\,E_\xi,0)$.
Subtracting the two translated values in \eqref{gt:curl}
therefore gives the exact identity
\begin{align}
 &\frac12\Big[
   (\operatorname{curl}_xf_d)_2(d+y_c,t)
   -(\operatorname{curl}_xf_d)_2(d-y_c,t)\Big]\notag\\
 &\hspace{12mm}
  =\frac{U_*\sqrt c\,e^{-c}}{L_*\tau_*}\sigma^{-2}
       \left[P_\beta(c)
         +\frac{\nu\tau_*}{L_*^2}
                  \sigma^{1-2\beta}Q(c)\right].
 \label{gt:opposite}
\end{align}
The source polynomials $P_\beta,Q$ are unchanged from
\eqref{eq:ag-polynomials}; neither the nonlinear term nor the
translation term was discarded before its exact evaluation
or cancellation.

For $\beta\ne1/6$ choose $c=c_*=(7-\sqrt{14})/2>0$.
The already verified identities
$Q(c_*)=0$ and
$P_\beta(c_*)=(1-6\beta)(10-4c_*)$ give the absolute value
of the right side of \eqref{gt:opposite} as
\[
 \frac{|U_*|\sqrt{c_*}\,e^{-c_*}}{L_*\tau_*}
    |1-6\beta|(10-4c_*)\,\sigma^{-2}.
\]
Its coefficient is strictly positive, since
$10-4c_*=-4+2\sqrt{14}>0$.
For $\beta=1/6$, choose $c=1$ instead. The right side is
\[
 \frac{U_*}{eL_*\tau_*}\sigma^{-2}
    \left[\frac{11}{3}
       +\frac{44\nu\tau_*}{L_*^2}\sigma^{2/3}\right].
\]
Its absolute value is at least
$11|U_*|\sigma^{-2}/(3eL_*\tau_*)$.
For any real numbers $a,b$,
$\max(|a|,|b|)\ge |a-b|/2$ by the triangle inequality.
Applying this to the two force-curl values proves that
their spatial supremum diverges in every case.

Finally, let any other pressure $\widetilde p_d$ and force
$\widetilde f_d$ solve the full equation with exactly the
same velocity $u_d$. Subtracting their equation from the
proved one gives
\[
 \widetilde f_d-f_d=\nabla_x(\widetilde p_d-p_d),\qquad
 \operatorname{curl}_x\widetilde f_d
       =\operatorname{curl}_xf_d.
\]
The second equality is also valid distributionally.
Every curl component is a difference of first spatial
derivatives of force components; uniform bounds on all
those derivatives would bound the curl. Equations
\eqref{gt:swirl}--\eqref{gt:opposite} contradict that bound.
This proves the assertion for every pressure.
\end{proof}

The target whole-space forcing condition
\eqref{eq:schwartz-force}, already at $K=0$ and first
spatial derivative order, requires the uniform bound
just excluded. The test locations may escape to infinity
when $d$ or $L$ does. The proof then concerns this global
condition; it does not claim a singularity at a fixed
finite spatial point. If the relevant test locations
stay in a fixed compact set, their divergent curl values
also preclude a smooth local extension there.
When $U_*=H_*=0$, every field in \eqref{gt:translation}
is zero for every such $d$ and $\beta$, so this degenerate
case has no obstruction.

There is a different moving-frame map which adds a
constant background velocity; its constant must also
be retained. For
$\widehat u(x,t)=u(x-d(t),t)+d'(t)$ and
$\widehat p(x,t)=p(x-d(t),t)$, the time derivative's
$-d'\cdot\nabla u$ term cancels the identical contribution
with positive sign in convection. The full resulting
force is $f(x-d,t)+d''(t)$. Equivalently the affine
pressure $\widehat p-d''(t)\cdot x$ removes this
acceleration from the force. Whenever $d'(t)\ne0$,
however, $\widehat u(t)$ has infinite energy:
if it and the translated Gaussian were both in $L^2$,
their difference, the nonzero constant vector $d'(t)$,
would be in $L^2$, whereas its squared integral over
a radius-$R$ ball is $|d'(t)|^2(4\pi R^3/3)$.
Thus this complete alternative map does not evade the
energy requirement. These results concern the specified
Gaussian fields and their exact moving-centre maps;
they do not exclude different profiles or dynamics.
